\begin{explanationbox}
	\textbf{\textcolor{CExplanation}{一句话直觉}}

	椭圆上的点到焦点和定直线的距离恒为正比例关系，比值即离心率 $e$。

	\medskip
	\textbf{\textcolor{CExplanation}{核心拆解}}
	\begin{description}[style=nextline, leftmargin=2.8em, labelsep=0.5em, itemsep=0.45em]
		\item[\textbf{要点一（必需三要素）：}] 定义严格依赖“一点、一焦点、一准线”的组合，缺一不可。
		\item[\textbf{要点二（离心率角色）：}] $e$ 既是比例常数，也是决定椭圆扁平程度的量（$e$ 越小越圆）。
	\end{description}

	\medskip
	\textbf{\textcolor{CExplanation}{几何本质（焦点准线协同）}}

	椭圆第二定义揭示了焦点和准线对曲线的双重约束，这种到定点与定直线的距离比例关系是圆锥曲线的统一框架（椭圆、双曲线、抛物线共用）。

	\medskip
	\textbf{\textcolor{CExplanation}{代数意义（坐标运算转化）}}

	定义式 $\frac{|PF|}{d(P,l)} = e$ 直接生成含根号的坐标方程，化简即得椭圆标准形式，是定义与方程互推的代数桥梁。

	\medskip
	\textbf{\textcolor{CExplanation}{考点价值}}
	\begin{itemize}[leftmargin=*, itemsep=0.5em, topsep=0.4em]
		\item \textbf{考法一（直接求曲线方程）：} 题干给出焦点、准线、$e$，直接套定义列方程。
		\item \textbf{考法二（巧解焦半径问题）：} 将 $|PF|$ 替换为 $e \cdot d(P,l)$，整体转化距离求解最值。
	\end{itemize}

	\medskip
	\textbf{\textcolor{CExplanation}{顿悟点}}

	\textbf{第二定义把“到两焦点的距离和”变为“到一焦点与一准线的距离比”，打开了全新的几何视角，是解题中转化的利器。}

	\medskip
	\textbf{\textcolor{CExplanation}{使用场景}}
	\begin{itemize}[leftmargin=*, itemsep=0.5em, topsep=0.4em]
		\item \textbf{场景一（定义法）：} 题目明确出现焦点和准线，求曲线方程或离心率。
		\item \textbf{场景二（转化法）：} 题目中焦半径系数与离心率 $e$ 明显匹配，需要转化距离关系。
	\end{itemize}
\end{explanationbox}
